% This document was generated by an LLM.

\documentclass{essay}

\title{Upper Bounds, Suprema, Maxima, and Sequence Limits}
\author{}
\date{}

\begin{document}

\maketitle

\section{Introduction}

Several concepts in elementary real analysis express, in different ways, the idea that a mathematical object does not extend indefinitely upward. Among the most important are the notions of an \emph{upper bound}, a \emph{supremum}, a \emph{maximum}, and the \emph{limit of a sequence}. At first sight these notions can appear to be variations on the same theme. An upper bound lies above a set, a supremum is the lowest possible such bound, a maximum is the largest element actually attained by the set, and the limit of an increasing sequence often seems to play the same role as its supremum.

There is indeed a close relationship among these ideas, but they belong to somewhat different parts of mathematics. Upper bounds, suprema, and maxima are primarily \emph{order-theoretic} notions: they concern the relative position of elements in an ordered set. A sequence limit is primarily a \emph{topological} or \emph{metric} notion: it concerns arbitrarily close approximation. In the real numbers these structures interact so well that the distinctions can be easy to overlook. In incomplete ordered fields such as $\mathbb{Q}$, however, the differences become much more visible.

The purpose of this essay is to compare the four notions precisely, to explain which implications hold between them, and to show how their relationship reveals one of the central ideas of real analysis: completeness.

\section{Upper Bounds}

Let $F$ be an ordered field and let $A\subseteq F$ be a nonempty set. An element $u\in F$ is called an \emph{upper bound} of $A$ if
\[
    a\leq u
    \qquad\text{for every }a\in A.
\]

The important point is that an upper bound need not belong to $A$. It is merely an element of the ambient ordered field that lies above every element of the set.

For example, if
\[
    A=(0,1)\subseteq\mathbb{R},
\]
then $1$ is an upper bound, but so are $2$, $100$, and every real number greater than or equal to $1$. Thus the existence of an upper bound is a relatively weak statement. It says only that the set is bounded from above.

It is useful to collect all upper bounds into a set:
\[
    U(A)=\{u\in F : a\leq u\text{ for every }a\in A\}.
\]
Then saying that $A$ is bounded above means simply that $U(A)$ is nonempty.

An upper bound therefore does not identify a distinguished endpoint. Usually there are many upper bounds. The notion records only that some ceiling exists.

\section{The Supremum}

The \emph{supremum}, or \emph{least upper bound}, is a refinement of the notion of upper bound. An element $s\in F$ is the supremum of $A$ if two conditions hold:
\begin{enumerate}
    \item $s$ is an upper bound of $A$;
    \item every upper bound $u$ of $A$ satisfies $s\leq u$.
\end{enumerate}

We write
\[
    s=\sup A.
\]

Equivalently, $\sup A$ is the minimum element of the set of all upper bounds:
\[
    \sup A=\min U(A),
\]
provided such a minimum exists.

The supremum is therefore much more specific than an arbitrary upper bound. If an upper bound merely says ``the set does not go above here,'' the supremum says ``this is the lowest possible ceiling.''

For example, for
\[
    A=(0,1)\subseteq\mathbb{R},
\]
we have
\[
    \sup A=1.
\]
Yet $1\notin A$. This already shows that a supremum need not be an element of the original set.

A very useful equivalent characterization is the approximation property. If $s=\sup A$, then for every $\varepsilon>0$ there exists some $a\in A$ such that
\[
    s-\varepsilon<a\leq s.
\]
Thus elements of $A$ can be found arbitrarily close to $s$ from below. If this were not true, then $s-\varepsilon$ would itself be an upper bound, contradicting the minimality of $s$.

This approximation property begins to resemble the language of limits. That resemblance is not accidental, but the two notions are still logically distinct.

\section{Maximum}

An element $m\in A$ is called the \emph{maximum} of $A$ if
\[
    a\leq m
    \qquad\text{for every }a\in A.
\]
We write
\[
    m=\max A.
\]

A maximum is therefore an upper bound that is actually attained by the set.

This gives an immediate implication:
\[
    \max A \text{ exists}
    \quad\Longrightarrow\quad
    \sup A \text{ exists and }\sup A=\max A.
\]

Indeed, if $m$ is the maximum, then $m$ is an upper bound. No smaller upper bound is possible, because $m\in A$, so every upper bound must be at least $m$.

The converse is false. A set can have a supremum without having a maximum. The standard example is
\[
    A=(0,1).
\]
Then
\[
    \sup A=1,
\]
but $A$ has no maximum because $1\notin A$.

Another useful example is the sequence range
\[
    A=\left\{1-\frac1n:n\in\mathbb{N}\right\}.
\]
Again,
\[
    \sup A=1,
\]
but no term of the set equals $1$, so there is no maximum.

Thus a maximum is, in an important sense, stronger than a supremum: it is not merely the least upper bound but an attained least upper bound.

The logical relation can be summarized as
\[
    \text{maximum} \Longrightarrow \text{supremum} \Longrightarrow \text{upper bound}.
\]
Neither implication reverses in general.

\section{Failure of Suprema in an Incomplete Field}

The distinction between upper bounds and suprema becomes especially important in $\mathbb{Q}$.

Consider
\[
    A=\{q\in\mathbb{Q}:q>0,\ q^2<2\}.
\]
The set is bounded above in $\mathbb{Q}$. For example, $2$ is an upper bound. In fact many rational numbers are upper bounds.

However, $A$ has no supremum in $\mathbb{Q}$. Informally, the least possible upper bound should be $\sqrt{2}$, but
\[
    \sqrt{2}\notin\mathbb{Q}.
\]
There is no rational number that fills the role.

Thus one can have
\[
    \text{upper bounds exist}
\]
without
\[
    \text{a supremum existing in the ambient field}.
\]

This is precisely the phenomenon ruled out by the least-upper-bound property of $\mathbb{R}$.

\section{Sequence Limits}

Now let $(a_n)$ be a sequence in an ordered field such as $\mathbb{R}$. We say that
\[
    a_n\to L
\]
if for every $\varepsilon>0$ there exists $N\in\mathbb{N}$ such that
\[
    n\geq N
    \quad\Longrightarrow\quad
    |a_n-L|<\varepsilon.
\]

This definition is fundamentally different from the definitions of upper bound, supremum, and maximum. It does not compare one element against every element of a set by the order relation. Instead, it concerns the eventual behavior of an indexed family of elements.

A sequence limit is therefore not inherently an upper bound. For example,
\[
    a_n=(-1)^n\frac1n
\]
converges to $0$, but $0$ is neither the maximum nor necessarily an upper bound in any conceptually important sense. The terms approach $0$ from both sides.

Similarly,
\[
    a_n=\frac1n
\]
converges to $0$, while every term satisfies $a_n>0$. The limit $0$ is not even an element of the range
\[
    \left\{\frac1n:n\in\mathbb{N}\right\}.
\]
Nor is it an upper bound of that range; it is a lower bound.

Thus in general there is no implication of the form
\[
    \text{sequence limit}\Longrightarrow\text{supremum}
\]
or vice versa.

The relation becomes much tighter when the sequence is monotone.

\section{Increasing Sequences: Where Limits and Suprema Meet}

Suppose $(a_n)$ is increasing:
\[
    a_1\leq a_2\leq a_3\leq\cdots.
\]
Let
\[
    A=\{a_n:n\in\mathbb{N}\}
\]
be its range.

If the sequence converges to $L$, then
\[
    L=\sup A.
\]

To see why, first note that $L$ must be an upper bound. If some term satisfied $a_k>L$, then monotonicity would give
\[
    a_n\geq a_k>L
\]
for every $n\geq k$, making convergence to $L$ impossible.

Next, no number smaller than $L$ can be an upper bound. If $u<L$, set
\[
    \varepsilon=L-u>0.
\]
Since $a_n\to L$, eventually
\[
    |a_n-L|<\varepsilon,
\]
which in particular gives
\[
    a_n>L-\varepsilon=u.
\]
Thus $u$ fails to be an upper bound.

Therefore $L$ is the least upper bound of the range.

Conversely, suppose the increasing sequence is bounded above and the supremum
\[
    s=\sup A
\]
exists. Then
\[
    a_n\to s.
\]
Indeed, given $\varepsilon>0$, the number $s-\varepsilon$ cannot be an upper bound of $A$. Hence some term satisfies
\[
    a_N>s-\varepsilon.
\]
Because the sequence is increasing, for every $n\geq N$,
\[
    s-\varepsilon<a_N\leq a_n\leq s,
\]
and therefore
\[
    |a_n-s|<\varepsilon.
\]

So for increasing sequences we obtain the precise identity
\[
    \boxed{\lim_{n\to\infty}a_n=\sup\{a_n:n\in\mathbb{N}\}}
\]
whenever either side exists under the appropriate hypotheses.

This formula is one of the clearest bridges between order theory and the theory of limits.

\section{A Hierarchy of Strength}

For a fixed set $A$, the three order-theoretic notions can be arranged naturally according to strength.

To say that $A$ has an upper bound means only that
\[
    U(A)\neq\varnothing.
\]

To say that $A$ has a supremum means that the set of upper bounds has a least element:
\[
    \min U(A)\text{ exists}.
\]

To say that $A$ has a maximum means that this least upper bound is also attained by $A$:
\[
    \sup A\in A.
\]

Hence
\[
    \boxed{
    \text{maximum}
    \Longrightarrow
    \text{supremum}
    \Longrightarrow
    \text{upper bound}
    }
\]
with strict separation between the notions.

The sequence limit does not fit into this same hierarchy because it concerns a different structure. A sequence may converge without being bounded above by its limit, and a set may have a supremum without having any sequence under discussion at all.

However, once one restricts attention to increasing sequences, the limit and supremum become equivalent descriptions of the same endpoint.

\section{Examples Separating the Four Notions}

It is useful to see examples in which the distinctions are explicit.

\subsection*{Upper bound without supremum}

In $\mathbb{Q}$, let
\[
    A=\{q\in\mathbb{Q}:q>0,\ q^2<2\}.
\]
Then $A$ has rational upper bounds but no supremum in $\mathbb{Q}$.

\subsection*{Supremum without maximum}

In $\mathbb{R}$, let
\[
    A=(0,1).
\]
Then
\[
    \sup A=1,
\]
but $A$ has no maximum.

\subsection*{Maximum}

Let
\[
    A=(0,1].
\]
Then
\[
    \max A=1,
\]
and consequently
\[
    \sup A=1.
\]

\subsection*{Limit unrelated to supremum}

For
\[
    a_n=\frac1n,
\]
we have
\[
    a_n\to0.
\]
But
\[
    \sup\left\{\frac1n:n\in\mathbb{N}\right\}=1.
\]
Thus the sequence limit and the supremum of its range are completely different.

\subsection*{Limit equal to supremum}

For
\[
    a_n=1-\frac1n,
\]
the sequence is increasing and
\[
    a_n\to1.
\]
Its range satisfies
\[
    \sup\left\{1-\frac1n:n\in\mathbb{N}\right\}=1,
\]
while no maximum exists.

This example displays three notions simultaneously:
\[
    \lim a_n=\sup A=1,
\]
but
\[
    \max A\text{ does not exist}.
\]

\section{The Role of the Ambient Space}

All four notions depend on the ambient mathematical structure.

For instance, consider a rational sequence converging in $\mathbb{R}$ to $\sqrt2$. One may construct an increasing sequence of rational numbers
\[
    a_n\in\mathbb{Q}
\]
with
\[
    a_n\to\sqrt2
\]
when the sequence is regarded as a sequence in $\mathbb{R}$.

But if the ambient metric space is $\mathbb{Q}$, then the same sequence does not converge, because a limit in $\mathbb{Q}$ must itself be rational.

Likewise, the set
\[
    A=\{q\in\mathbb{Q}:q>0,\ q^2<2\}
\]
has no supremum when regarded as a subset of $\mathbb{Q}$, but when regarded as a subset of $\mathbb{R}$ it has
\[
    \sup A=\sqrt2.
\]

Therefore statements such as ``$A$ has a supremum'' or ``$(a_n)$ converges'' are incomplete unless the ambient space is understood.

This is one reason completeness is a property of a space, not merely of a particular subset or sequence.

\section{Completeness as the Missing Link}

The real numbers satisfy the least-upper-bound property:

\medskip
\noindent
Every nonempty subset of $\mathbb{R}$ that is bounded above has a supremum in $\mathbb{R}$.
\medskip

This property immediately implies the monotone convergence theorem for sequences:

\medskip
\noindent
Every increasing sequence of real numbers that is bounded above converges to a real number.
\medskip

The proof works because if
\[
    A=\{a_n:n\in\mathbb{N}\},
\]
then completeness supplies
\[
    s=\sup A,
\]
and monotonicity forces
\[
    a_n\to s.
\]

In $\mathbb{Q}$ this argument can fail. A bounded increasing rational sequence may approach $\sqrt2$, but the corresponding range has no supremum in $\mathbb{Q}$, and the sequence has no rational limit.

Thus the failure of a supremum and the failure of a sequence limit are not two unrelated defects. They are two manifestations of the same incompleteness.

This is one of the central structural insights of real analysis. Completeness may be formulated in several apparently different languages: least upper bounds, convergence of Cauchy sequences, convergence of bounded monotone sequences, nested intervals, and others. Their equivalence reflects the fact that they are different ways of saying that the real number line has no gaps.

\section{Order-Theoretic Versus Topological Information}

It is useful to distinguish exactly what kind of information each notion contains.

An upper bound gives \emph{global order information}. It compares one element with every element of the set.

A supremum gives more refined global order information. It identifies the precise boundary between the set and its upper bounds.

A maximum gives both boundary information and attainment: the boundary point belongs to the set itself.

A sequence limit gives \emph{asymptotic approximation information}. It says that after discarding finitely many terms, every remaining term lies arbitrarily close to a distinguished point.

These are different logical forms. For an upper bound one writes
\[
    \forall a\in A,\quad a\leq u.
\]
For convergence one writes
\[
    \forall\varepsilon>0\ \exists N\ \forall n\geq N,
    \quad |a_n-L|<\varepsilon.
\]

The quantifier patterns themselves show the distinction. Bounding concerns all elements simultaneously. Convergence concerns eventual behavior at every prescribed scale of accuracy.

For monotone sequences, these two kinds of information reinforce each other. Order prevents oscillation, while the supremum provides the unique endpoint toward which the sequence must move.

\section{Conclusion}

Upper bounds, suprema, maxima, and sequence limits all express forms of ``endpoint'' behavior, but they do so in different ways.

An upper bound is merely a ceiling. A supremum is the lowest possible ceiling. A maximum is a ceiling that is actually attained. These three notions form a genuine hierarchy:
\[
    \text{maximum}\Longrightarrow\text{supremum}\Longrightarrow\text{upper bound}.
\]

A sequence limit is not part of this hierarchy. It belongs to the theory of approximation rather than directly to the theory of order. A general convergent sequence need not approach the supremum of its range, and a set with a supremum need not arise from any particular sequence.

The special case of monotone sequences brings the two worlds together. For an increasing sequence,
\[
    \lim a_n=\sup\{a_n:n\in\mathbb{N}\},
\]
whenever the relevant endpoint exists. This identity is not merely a useful theorem. It exposes the connection between order and convergence that lies at the heart of completeness.

The contrast between $\mathbb{Q}$ and $\mathbb{R}$ makes the point especially clear. In $\mathbb{Q}$, a set may be bounded above without possessing a supremum, and a bounded increasing sequence may fail to converge. In $\mathbb{R}$, completeness fills these gaps. Suprema exist where order demands them, and monotone sequences converge toward those suprema.

The four notions are therefore best understood not as competing definitions of one vague idea of a ``limit,'' but as distinct concepts whose interaction reveals the structure of the real number system.

\end{document}
